\section{Occupation measures and the whole trace}
\label{sec:occupation}

Proposition~\ref{prop:disk-calibration} compares a microscopic segment
with a disk whose radius is a fixed fraction of the domain size.  We
first make this comparison inside second moments.  The resulting
approximation permits us to use ordered-curve convergence at a fixed
disk radius, and only afterwards let that radius tend to zero.  We
also control the boundary strip, since weak convergence of finite
measures on the closed square includes convergence of total mass.

\subsection{Segment probes and their replacement by disks}

Let $\mathcal E_n$ be the deterministic set of full medial segments
available in the discrete square, including the boundary collar of
Section~\ref{sec:introduction}.  A segment here is an edge of the drawn
medial graph, before applying $\Pi$.  Write $z_e$ for its projected
midpoint and $I_e$ for the indicator that it belongs to the distinguished
strand.  The pairing at every medial vertex splits that vertex into
two vertices of degree two.  In the resulting finite graph, the
component containing a distinguished terminal half-edge is a path,
with the other distinguished half-edge as its other end.  It cannot
repeat an edge: a repeated edge would put that component on a cycle,
which cannot contain its degree-one terminal vertices.  Consequently
the counting multiplicity in~\eqref{eq:counting-measure} is exactly
$I_e$, even when the unsplit medial vertex is visited twice.  Distinct
collar edges are still counted separately if their projections coincide.

There is a constant $A_0>0$, determined by the density of medial
midpoints in the stated lattice units, such that
\begin{equation}
 A_0n^{-2}\sum_{e\in\mathcal E_n}h(z_e)
 \longrightarrow \int_D h(z)\dd A(z)
 \qquad\bigl(h\in C(\overline D)\bigr).
 \label{eq:probe-riemann}
\end{equation}
Indeed the bulk midpoints form a periodic set of fixed density, and
there are only $O(n)$ additional boundary probes.  Define
\begin{equation}
 X_n(f)=A_0n^{-2}\sum_{e\in\mathcal E_n}f(z_e)n^x I_e
       =A_0n^{-d}\sum_{e\in\mathcal E_n}f(z_e)I_e.
 \label{eq:microscopic-test}
\end{equation}

Use the sequence of radii $\epsilon\downarrow0$ supplied by
Proposition~\ref{prop:disk-calibration}, together with its constants
$L_\epsilon\asymp\epsilon^{-x}$.  The physical scale of a unit tile is
$a_0/n$.  For a probe $e$ in the bulk, let $B_e^\epsilon$ be the
whole-tile disk of radius $s=\epsilon n$ used in that proposition,
translated to the probe.  Its physical radius is $a_0\epsilon$, up to
$O(n^{-1})$.  Let $J_e^\epsilon$ indicate that the distinguished strand
enters this disk; the occurrence of an unrelated loop does not count.
For $f\in C_c(D)$, put
\begin{equation}
 X_{n,\epsilon}(f)
   =A_0n^{-2}\sum_{e\in\mathcal E_n}f(z_e)L_\epsilon J_e^\epsilon.
 \label{eq:disk-test}
\end{equation}
For fixed $\epsilon$ sufficiently small, every disk occurring in this
sum is inside $D$ once $n$ is sufficiently large.

\begin{proposition}[Replacement in second moment]
\label{prop:l2-replacement}
For every real $f\in C_c(D)$,
\begin{equation}
 \lim_{\epsilon\downarrow0}\limsup_{n\to\infty}
 \E\bigl[|X_n(f)-X_{n,\epsilon}(f)|^2\bigr]=0,
 \label{eq:l2-replacement}
\end{equation}
where $\epsilon$ tends to zero through the sequence in
Proposition~\ref{prop:disk-calibration}.  The same approximation holds
for every finite vector of such tests.
\end{proposition}

\begin{proof}
Fix a compact set $K\Subset D$ containing $\supp f$, and abbreviate
\[
 H_e^0=I_e,\qquad H_e^1=J_e^\epsilon,
 \qquad w_0=n^x,\qquad w_1=L_\epsilon,
 \qquad V_e^a=w_aH_e^a\quad(a\in\{0,1\}).
\]
We first compare products at two probes $e,e'$ satisfying
$z_e,z_{e'}\in K$ and $|z_e-z_{e'}|\ge g$, where $g>0$ is fixed.
Constants in this comparison may depend on $K$ and $g$.  Choose
$\ell=c_{K,g}n$ in tile units, with $c_{K,g}>0$ small enough that
balls of radius several times $\ell$ about the two probes are disjoint
and contained in the square.  All auxiliary cuts below leave fixed
annular collars.  Let $\mathcal F$ be the edges outside the ball of
radius $\ell$ about $e$.

Take $R\asymp\epsilon^b n$ as in
Proposition~\ref{prop:disk-calibration}, where $b>0$ will be sufficiently
small, and write $h_\epsilon=R/n\asymp\epsilon^b$.
Let $S_R^a$ denote the local two-color reaching event from the segment
start when $a=0$ and from $B_e^\epsilon$ when $a=1$.  For sufficiently
small $\epsilon$, these events and their collars are inside the
$\ell$-ball.  On every exterior configuration of positive probability,
set
\[
 p_a^{\mathcal F}=\Pp_n(S_R^a\mid\mathcal F),
 \qquad
 Q_a^{\mathcal F}=\Pp_n(\,\cdot\mid\mathcal F,S_R^a),
 \qquad \Lambda_a=w_ap_a^{\mathcal F}.
\]
The conditional FK law has the partition induced by the exposed
exterior; the uniformity in that partition is essential here.

Write $p_a$ for the corresponding ordinary plane probability.  The
relative comparison through the full collar in Input~\ref{inp:fk}
gives, for some $c_{\rm mix}>0$,
\begin{equation}
 p_a^{\mathcal F}=p_a
       \bigl(1+O_{K,g}(h_\epsilon^{c_{\rm mix}})\bigr).
 \label{eq:occupation-relative-comparison}
\end{equation}
Proposition~\ref{prop:sharp-arms}, Equation~\eqref{eq:start-probabilities},
and Proposition~\ref{prop:disk-calibration} consequently imply
\begin{align}
 \Lambda_a&\le C_{K,g}h_\epsilon^{-x},
 \label{eq:occupation-lambda-bound}\\
 |\Lambda_0-\Lambda_1|
 &\le C_{K,g}h_\epsilon^{-x}
 \bigl(\epsilon^{c_{\rm cal}}
          +h_\epsilon^{c_{\rm mix}}+o_{n,\epsilon}(1)\bigr)
 \label{eq:occupation-lambda-difference}
\end{align}
for some $c_{\rm cal}>0$.  Here and below $o_{n,\epsilon}(1)$ tends to
zero as $n\to\infty$ for each fixed admissible $\epsilon$; it can be
chosen uniformly over the probes and exposed exteriors under
consideration.  Translation, rotation, and the color exchange in
Proposition~\ref{prop:disk-calibration} identify their plane ratios.

Fix a second probe type $v\in\{0,1\}$.  Let $T_{e'}^v$ be its
two-color reaching event to a cut of radius $\ell/C$ about $e'$, with
$C$ a sufficiently large fixed constant, defined using actual paths
inside that cut rather than connections supplied by exterior wiring.
This event is
$\mathcal F$-measurable and is necessary for $H_{e'}^v=1$.
Define
\[
 q_a^{\mathcal F}
   =\E_{Q_a^{\mathcal F}}[H_e^aH_{e'}^v].
\]
Proposition~\ref{prop:outward} couples $Q_0^{\mathcal F}$ and
$Q_1^{\mathcal F}$ with failure probability at most
$C(s/R)^{c_{\rm out}}$, for some $c_{\rm out}>0$.  On success there is
an identical intermediate band with exactly two through strands, and
the completions outside that band agree.  The two through strands have
one possible connection through the interior.  In the segment
experiment that connection contains $e$, while in the disk experiment
it enters $B_e^\epsilon$.  Whether the connection belongs to the
distinguished strand is therefore the same in both completions.
The identical outer connection also preserves its hit verdict at $e'$.
This last observation is needed because $H_{e'}^v$ itself need not be
$\mathcal F$-measurable.  It gives
\begin{equation}
 |q_0^{\mathcal F}-q_1^{\mathcal F}|
 \le C_{K,g}(s/R)^{c_{\rm out}}\1_{T_{e'}^v}.
 \label{eq:occupation-q-difference}
\end{equation}

A hit by the distinguished strand also requires continuation from
$R$ to a cut of radius comparable to $\ell$ about $e$.  By the start
comparisons, sharp arms, and their conditional collar versions,
\begin{equation}
 q_a^{\mathcal F}
    \le C_{K,g}h_\epsilon^x\1_{T_{e'}^v}.
 \label{eq:occupation-q-bound}
\end{equation}
One can obtain this bound by taking the ratio of the start probability
to radius $\ell/C$ and that to radius $R$ in the conditional FK law.
Both cuts have collars before the exposed exterior; the start factors
in Equation~\eqref{eq:start-probabilities} cancel, leaving
$\pi_2(R,\ell/C)\le C_{K,g}h_\epsilon^x$.

Since $H_e^a=1$ implies $S_R^a$, the conditional expectation of
$V_e^aH_{e'}^v$ is $\Lambda_aq_a^{\mathcal F}$.  Subtracting these
expressions and using
Equations~\eqref{eq:occupation-lambda-bound}--\eqref{eq:occupation-q-bound}
yields
\begin{align}
 &\left|\E\bigl[(V_e^0-V_e^1)H_{e'}^v\mid\mathcal F\bigr]\right|
 \nonumber\\
 &\quad\le C_{K,g}\1_{T_{e'}^v}
 \left[h_\epsilon^{-x}(s/R)^{c_{\rm out}}
       +\epsilon^{c_{\rm cal}}+h_\epsilon^{c_{\rm mix}}
       +o_{n,\epsilon}(1)\right].
 \label{eq:conditional-product-replacement}
\end{align}
Thus the continuation bound cancels the factor $h_\epsilon^{-x}$ in
the normalization error.  The coupling error retains this factor.
Finally, the necessary arms at the second probe satisfy
\[
 w_v\Pp_n(T_{e'}^v)\le C_{K,g}.
\]
Multiplication of Equation~\eqref{eq:conditional-product-replacement}
by $w_v$ and integration therefore show that
\begin{equation}
 \limsup_{n\to\infty}
 \sup_{\substack{z_e,z_{e'}\in K\\|z_e-z_{e'}|\ge g}}
 \left|\E[(V_e^0-V_e^1)V_{e'}^v]\right|
 \le C_{K,g}\eta(\epsilon),
 \label{eq:separated-product-replacement}
\end{equation}
where
\[
 \eta(\epsilon)=\epsilon^{(1-b)c_{\rm out}-bx}
                +\epsilon^{c_{\rm cal}}
                +\epsilon^{bc_{\rm mix}}.
\]
Choose $b$ sufficiently small that
$(1-b)c_{\rm out}-bx>0$, in addition to the requirements in
Proposition~\ref{prop:disk-calibration}.  Then
$\eta(\epsilon)\to0$.  Interchanging the probes gives the analogous
replacement in the second factor.

It remains to control nearby probes.  The first-moment arm bound on
$K$ is
\begin{equation}
 \E V_e^a\le C_K\qquad(a=0,1).
 \label{eq:interior-probe-first-moment}
\end{equation}
For two probes at distance $y\ge C\epsilon$, impose their necessary
arms separately to radii comparable to $y$, using disjoint collars.
Conditioning on the exterior of either collar leaves the other event
measurable, so the uniform conditional bounds multiply.  They give
\begin{equation}
 \E[V_e^aV_{e'}^v]\le C_K y^{-2x}.
 \label{eq:probe-pair-bound}
\end{equation}
When $a=v=0$, this holds down to a fixed multiple of the mesh; at
smaller distances use the bound $n^{2x}$.  Thus its right side can
there be replaced by $C_K(y\vee n^{-1})^{-2x}$.
If $y\le C\epsilon$ and at least one probe is a disk probe,
Equation~\eqref{eq:interior-probe-first-moment} instead gives
\begin{equation}
 \E[V_e^aV_{e'}^v]\le C_K\epsilon^{-x}.
 \label{eq:overlapping-disk-bound}
\end{equation}
For a mixed pair bound the disk factor by $L_\epsilon$ and use the
microscopic first moment; for two disk probes use the first moment of
either one.

The periodic probe set has at most $C(1+n^2u^2)$ points in a ball of
radius $u$.  Summing Equation~\eqref{eq:probe-pair-bound} in dyadic
distance annuli with $y<g$ shows that its contribution to the double
Riemann sum is at most
$C_K(g^{2-2x}+n^{2x-2})$.  The contribution of
Equation~\eqref{eq:overlapping-disk-bound}, after taking the upper
limit in $n$, is at most $C_K\epsilon^{2-x}$.  Both exponents are
positive.  Expanding the square in
Equation~\eqref{eq:l2-replacement}, use these bounds for pairs at
distance less than $g$ and
Equation~\eqref{eq:separated-product-replacement} for all other pairs.
For fixed $g$ and sufficiently small $\epsilon$ this proves
\[
 \limsup_n\E|X_n(f)-X_{n,\epsilon}(f)|^2
 \le C_K\|f\|_\infty^2
       \bigl(g^{2-2x}+\epsilon^{2-x}\bigr)
       +C_{K,g}\|f\|_\infty^2\eta(\epsilon).
\]
Let $\epsilon\downarrow0$ first and then $g\downarrow0$.
Applying the scalar result to each component proves the assertion for
finite vectors.
\end{proof}

\subsection{Minkowski content in the closed square}

We record the whole-domain consequence of the SLE content theorems,
including the boundary estimate needed to match the convention of
Equation~\eqref{eq:content-measure}.

\begin{lemma}[Whole-square content]
\label{lem:whole-content}
Let $\eta$ be chordal $\mathrm{SLE}_\kappa$ in $(D;a,b)$, with the
parameters of Equation~\eqref{eq:parameters}.  Its interior
$d$-dimensional Euclidean Minkowski-content measure extends by zero
on $\partial D$ to a finite random measure $\mu_\eta$ on
$\overline D$.  For every real $f\in C(\overline D)$,
\begin{equation}
 r^{-x}\int_\C f(\Pi z)
       \1_{\{\dist(z,\tr\eta)<r\}}\dd A(z)
 \longrightarrow \mu_\eta(f)
 \quad\hbox{in probability as }r\downarrow0.
 \label{eq:whole-content-convergence}
\end{equation}
Its expected total mass is finite, and some nonnegative
$f\in C_c(D)$ has $\E\mu_\eta(f)>0$.  If
$U_\delta=\{z\in\overline D:\dist(z,\partial D)\le\delta\}$,
then
\begin{equation}
 \E\mu_\eta(U_\delta)\le C\delta^{1-x}.
 \label{eq:continuum-boundary-strip}
\end{equation}
\end{lemma}

\begin{proof}
Here $4<\kappa\le6<8$ and $0<x<1$.  Theorem~3.1 of
\cite{LawlerRezaeiContent} establishes almost sure and $L^2$
convergence of the full-trace neighborhood-area approximations on
interior dyadic squares in the half-plane.  Together with its
measure construction and Theorem~1.1, this gives the interior content
measure and convergence against continuous tests of compact support.
The normalization there is precisely radius to the power $d-2$
times Euclidean neighborhood area.

For completeness, this interior conclusion is preserved by the
conformal map $F:\mathbb H\to D$ taking $0,\infty$ to $a,b$.
The inverse image of a compact subset of $D$ is a compact subset of
$\mathbb H$.  Given $\vartheta>0$, subdivide a neighborhood of that
compact set into patches with centers $w_j$ on which
$|F'|/|F'(w_j)|$ lies between $1-\vartheta$ and $1+\vartheta$.
For sufficiently small $r$, conformal distortion bounds the preimage
of a physical $r$-neighborhood on the $j$th patch between neighborhoods
with radii $(1\pm C\vartheta+o_r(1))r/|F'(w_j)|$.
Neighborhoods of the whole
trace may be used: a trace portion outside a larger patch cannot
affect a sufficiently small neighborhood of its inner patch.
Changing area by $|F'|^2$ and applying the half-plane convergence
squeezes the limit between the corresponding patchwise integrals.
Taking $r\downarrow0$ first and then $\vartheta\downarrow0$ identifies
the measure as
\[
 F_*\bigl(|F'|^d\mu_{\mathbb H}\bigr)
 \quad\hbox{on }D,
\]
and gives convergence in probability for compactly supported tests.
Continuous nonnegative tests can first be approximated on these
patches; the signed case follows by subtraction.

Theorem~2.3 of \cite{LawlerRezaeiBasic} applies to chordal SLE in any
simply connected domain.  It gives a Euclidean one-point asymptotic,
with an error bounded by a power of
$r/\dist(z,\partial D)$ when this ratio is below $1/10$.
Its Green function is a constant times the conformal radius to the
power $-x$, multiplied by the angular factor
$\sin^{8/\kappa-1}\theta$.  This factor is at most one, and Koebe's
theorem compares conformal radius to boundary distance.  Hence, with
$y=\dist(z,\partial D)$,
\begin{equation}
 \Pp\{\dist(z,\tr\eta)<r\}
       \le C(r/y)^x\qquad(0<r<y/10).
 \label{eq:sle-distance-upper-bound}
\end{equation}
The constants are uniform in $z\in D$, including near the corners.
The interior Green moment identity in
Theorem~3.1 of \cite{LawlerRezaeiContent}, transported by $F$, implies
positive expected content against every nonzero nonnegative
continuous interior test.  This uses the Green function with the
Euclidean content constant absorbed, as in that paper.

Let $M_r$ denote the random approximation measure on the left side
of Equation~\eqref{eq:whole-content-convergence}.  For $r<\delta/10$,
integrating Equation~\eqref{eq:sle-distance-upper-bound} over
$10r<y\le\delta$ gives
\[
 \E M_r\bigl(\{z\in D:10r<\dist(z,\partial D)\le\delta\}\bigr)
 \le C\int_0^\delta y^{-x}\dd y
 \le C\delta^{1-x}.
\]
The interior strip $y\le10r$ has area $O(r)$, so its contribution is
at most $Cr^{1-x}$ even without a probability bound.  Outside $D$,
only the $r$-neighborhood of $\overline D$ contributes.  Its area is
$O(r)$, and $\Pi$ maps it onto the boundary.  We have therefore shown
\begin{equation}
 \E M_r(U_\delta)
       \le C\delta^{1-x}+Cr^{1-x}.
 \label{eq:content-approximation-boundary}
\end{equation}

The interior Green bound is integrable over the square.  Local
convergence, Fatou's lemma, and an exhaustion of $D$ consequently
give a finite interior measure with finite expected total mass.
The same argument in boundary strips gives
Equation~\eqref{eq:continuum-boundary-strip}.  Extend this measure
by assigning zero mass to $\partial D$.  For a continuous test on
$\overline D$, multiply it by a continuous cutoff vanishing on
$U_\delta$ and equal to one outside $U_{2\delta}$.  Local convergence
applies to the cutoff test.  Equations~\eqref{eq:continuum-boundary-strip}
and~\eqref{eq:content-approximation-boundary}, followed by Markov's
inequality, make the discarded parts arbitrarily small, first as
$r\downarrow0$ and then as $\delta\downarrow0$.
This proves Equation~\eqref{eq:whole-content-convergence} with exactly
the prescribed projection and normalization.  The conformal map
extends to the marked corners by Jordan-domain boundary continuity,
so this argument includes both chordal endpoints.
\end{proof}

\subsection{Identification of the joint limit}

\begin{proof}[Proof of Theorem~\ref{thm:main}]
Fix an admissible $\epsilon>0$ and $f\in C_c(D)$.  Input~\ref{inp:trace}
implies, jointly with the ordered curve, that
\begin{equation}
 X_{n,\epsilon}(f)\ \Longrightarrow\
 L_\epsilon\int_D f(z)
       \1_{\{\dist(z,\tr\eta)<a_0\epsilon\}}\dd A(z).
 \label{eq:fixed-radius-joint-limit}
\end{equation}
To verify the approximation of the disk tests, consider any sequence
of curves converging in the ordered-curve metric.  Their trace sets
then converge in Hausdorff distance.  The tile
rounding, disk center choice, and rounded switch drawing change
distances by $O(n^{-1})$.  Sandwich the disk indicators between
Euclidean distance indicators with radii
$a_0\epsilon\pm o(1)$ and apply
Equation~\eqref{eq:probe-riemann}.
For every closed set $T\subset\C$ and every $r>0$, the level set
$\{z:\dist(z,T)=r\}$ has area zero: the distance function has gradient
of magnitude one almost everywhere off $T$, whereas its gradient
vanishes almost everywhere on any positive-area level set.
Thus the sandwich has the asserted limit, and the continuous mapping
theorem gives Equation~\eqref{eq:fixed-radius-joint-limit}.
This reasoning uses no simplicity of $\tr\eta$ and applies to any
finite collection of tests.

The deterministic quantities $L_\epsilon(a_0\epsilon)^x$ are bounded
above and below by positive constants.  Choose a subsequence of the
admissible radii, still denoted by $\epsilon$, on which
\begin{equation}
 L_\epsilon(a_0\epsilon)^x\longrightarrow L\in(0,\infty).
 \label{eq:occupation-limit-constant}
\end{equation}
By Lemma~\ref{lem:whole-content}, the right side of
Equation~\eqref{eq:fixed-radius-joint-limit} tends in probability to
$L\mu_\eta(f)$ on the same SLE probability space.  The finite-vector
version of Proposition~\ref{prop:l2-replacement} and the
converging-together argument now give
\begin{equation}
 \bigl([\eta_n],X_n(f_1),\ldots,X_n(f_k)\bigr)
 \ \Longrightarrow\
 \bigl([\eta],L\mu_\eta(f_1),\ldots,L\mu_\eta(f_k)\bigr)
 \label{eq:interior-joint-limit}
\end{equation}
for all real $f_1,\ldots,f_k\in C_c(D)$, along all integers $n$.
Indeed, at each fixed radius the approximation converges along the
entire mesh sequence; its replacement error tends to zero in the
successive order $n\to\infty$, $\epsilon\downarrow0$.
Extracting radii in Equation~\eqref{eq:occupation-limit-constant}
does not extract meshes.  This also shows directly why the limiting
measure and curve in Equation~\eqref{eq:interior-joint-limit} belong
to the same SLE sample.

To pass to the closed square, regard $X_n$ as the measure defined by
Equation~\eqref{eq:microscopic-test}.  If a probe is at physical
boundary distance $y\ge C/n$, a hit requires its two color arms to a
cut of radius comparable to $ny$ in tile units.  That cut and a full
collar can be placed inside the square.  The conditional arm bound
from Proposition~\ref{prop:sharp-arms} gives
\begin{equation}
 \E[n^x I_e]\le Cy^{-x}.
 \label{eq:lattice-near-boundary-first-moment}
\end{equation}
There are $O(n)$ probes in every strip of width $1/n$ adjacent to the
square boundary.  Consequently, for a sufficiently large fixed
integer $j_0$ and $\delta\ge j_0/n$,
\begin{align}
 \E X_n(U_\delta)
 &\le Cn^{x-1}
       +Cn^{-1}\sum_{j=j_0}^{\lceil\delta n\rceil}(j/n)^{-x}
 \nonumber\\
 &\le Cn^{x-1}+C\delta^{1-x}.
 \label{eq:lattice-boundary-strip}
\end{align}
The first term includes all probes in the remaining bounded-width
mesh strip and in the collar, using $I_e\le1$.  Projecting them does
not change their number.  The terminal connectors have no mass.
Thus all boundary contributions vanish in the successive limits
$n\to\infty$, $\delta\downarrow0$.

Use the same continuous interior cutoffs as in the proof of
Lemma~\ref{lem:whole-content}.  For every $f\in C(\overline D)$ the
discarded discrete part has first moment at most
$C\|f\|_\infty(\delta^{1-x}+n^{x-1})$, and the discarded limit has
first moment at most $CL\|f\|_\infty\delta^{1-x}$.
Equation~\eqref{eq:interior-joint-limit} therefore extends to every
finite set of continuous tests on $\overline D$, including $f=1$.
The same first-moment estimates give
$\sup_n\E X_n(\overline D)<\infty$, after absorbing finitely many
small meshes.  For each finite $M$, the set of nonnegative measures of
total mass at most $M$ is compact in the weak topology on the compact square,
so the measures are tight.  Together with curve tightness from
Input~\ref{inp:trace}, a countable determining family of continuous
tests identifies the joint limit as
\[
 ([\eta_n],X_n)\ \Longrightarrow\ ([\eta],L\mu_\eta).
\]
Finally, by the edge-counting identity preceding
Equation~\eqref{eq:microscopic-test},
$X_n=A_0n^{-d}\sum_{j=1}^{N_n}\delta_{\Pi(z_{n,j})}$.
Taking the single deterministic constant
\[
 c(q)=\frac{A_0}{L}\in(0,\infty)
\]
proves the asserted convergence in the product of the ordered-curve
and weak finite-measure topologies, including convergence of total
mass.
\end{proof}
